
% Paper E: CC-Image Vertex Census Sequence
\documentclass[11pt,a4paper]{article}
\usepackage[utf8]{inputenc}
\usepackage{amsmath,amssymb,amsthm}
\usepackage{geometry}
\usepackage{booktabs}
\usepackage{array}
\usepackage{hyperref}

\geometry{margin=2.5cm}

\theoremstyle{definition}
\newtheorem{definition}{Definition}[section]
\newtheorem{theorem}{Theorem}[section]
\newtheorem{proposition}{Proposition}[section]
\newtheorem{lemma}{Lemma}[section]
\newtheorem{corollary}{Corollary}[section]
\newtheorem{remark}{Remark}[section]

\title{OEIS Sequence Proposal: CC-Image Vertex Census\\
\large Counting the Graph-Theoretic Address Space of the Clique-Cluster Transformation}
\author{Alejandro Zarzuelo Urdiales\\
Department of Mathematics, University of Oviedo\\
UO275823@uniovi.es}
\date{May 2026}

\begin{document}
\maketitle

\begin{abstract}
This paper proposes a new integer sequence for the OEIS that counts, for each $N \geq 2$, the number of non-isomorphic simple graphs on exactly $N$ vertices that arise as CC-images of complete $p$-multigraphs under the Clique-Cluster Transformation. The resulting sequence is related to divisor-counting functions but differs fundamentally at pronic numbers $N = m(m+1)$, where a structural bonus accounts for the ``phase transition'' in the coding clique size formula. We derive a closed-form expression, compute values for $N \leq 60$, and establish connections to OEIS A038548, A056924, and A002378.
\end{abstract}

\section{Introduction}

The Clique-Cluster Transformation $f_{CC}$, introduced in \cite{zarzuelo2025}, encodes $p$-multigraphs into simple graphs by replacing each vertex with a coding clique of size $K$ and using index-rigid inter-cluster wiring to encode edge multiplicities. A fundamental parameter is the vertex count of the CC-image: for a $p$-multigraph $G$ on $n$ vertices, the CC-image has exactly $n \cdot K$ vertices, where $K = \max(P, \omega(G)) + 1$ under the Clique-Size Dominance condition.

This paper addresses the natural enumeration question: \emph{For a given number of vertices $N$, how many distinct CC-images of complete $p$-multigraphs exist?} This counts the ``address space'' of the CC-transformation --- the number of distinct positions in the universe of simple graphs that the encoding can occupy.

\section{Mathematical Framework}

\subsection{Vertex Count of CC-Images}

\begin{definition}[CC-Image Vertex Count]
For a complete $p$-multigraph $K_n^{(p)}$ on $n$ vertices, the CC-image under the minimum valid coding clique size has:
\[
K = \max(p, n) + 1,
\]
and therefore the CC-image has:
\[
N = n \cdot K = n \cdot (\max(p, n) + 1)
\]
vertices.
\end{definition}

The piecewise structure of $\max(p, n)$ creates two regimes:

\begin{proposition}[Regime Decomposition]
\label{prop:regime}
For $n \leq p$: $N = n(p + 1)$. \quad For $n > p$: $N = n(n + 1)$.
\end{proposition}

\subsection{The Proposed Sequence}

\begin{definition}[Sequence E: CC-Image Vertex Census]
For $N \geq 2$, the sequence $a(N)$ counts the number of pairs $(n, p)$ with $n \geq 1$, $p \geq 1$ such that:
\[
N = n \cdot (\max(p, n) + 1).
\]
Equivalently, $a(N)$ is the number of non-isomorphic simple graphs on exactly $N$ vertices that are CC-images of complete $p$-multigraphs.
\end{definition}

\begin{remark}[Equivalence of Counting]
Each valid pair $(n, p)$ produces a unique (up to isomorphism) CC-image because the cluster structure --- $n$ coding cliques of size $K$ with all-to-all inter-cluster edges of multiplicity $p$ --- uniquely determines $(n, p)$ from the isomorphism type of the resulting simple graph. Specifically:
\begin{itemize}
\item The size $K$ of the maximal coding cliques determines $\max(p, n) = K - 1$.
\item The number $n$ of such cliques determines $n$.
\item The number of inter-cluster edges per pair of clusters determines $p$.
\end{itemize}
Therefore, counting pairs $(n, p)$ is equivalent to counting non-isomorphic CC-images.
\end{remark}

\section{Main Results}

\subsection{Closed-Form Expression}

\begin{theorem}[Formula for $a(N)$]
\label{thm:formula}
For $N \geq 2$:
\[
a(N) = \left| \left\{ n \mid N \;\middle|\; \frac{N}{n} \geq n + 1 \right\} \right| + \begin{cases} m - 1 & \text{if } N = m(m+1) \text{ for some } m \geq 2, \\ 0 & \text{otherwise.} \end{cases}
\]
\end{theorem}

\begin{proof}
We count pairs $(n, p)$ satisfying $N = n \cdot (\max(p, n) + 1)$ in two cases.

\textbf{Case 1: $n \leq p$.} Then $N = n(p+1)$, so $p = N/n - 1$. The conditions $n \geq 1$, $p \geq 1$, and $n \leq p$ require:
\begin{itemize}
\item $n \mid N$,
\item $N/n - 1 \geq 1$, i.e., $N/n \geq 2$, i.e., $n \leq N/2$,
\item $N/n - 1 \geq n$, i.e., $N \geq n(n+1)$.
\end{itemize}
The binding constraint is $N/n - 1 \geq n$, i.e., $N \geq n(n+1)$. The number of such divisors is:
\[
c_1(N) = \left| \left\{ n \mid N \;\middle|\; \frac{N}{n} \geq n + 1 \right\} \right|.
\]

\textbf{Case 2: $n > p$.} Then $N = n(n+1)$, so $N$ must be a pronic (oblong) number. If $N = m(m+1)$ for some $m$, then $n = m$ and $p$ ranges over $\{1, 2, \ldots, m-1\}$ (since $p < n = m$). The number of such pairs is $m - 1$.

Combining: $a(N) = c_1(N) + c_2(N)$, where $c_2(N) = m - 1$ if $N = m(m+1)$, and $0$ otherwise.
\end{proof}

\subsection{Relationship to Divisor-Counting Sequences}

\begin{proposition}[Connection to A056924]
\label{prop:A056924}
The first term of Theorem \ref{thm:formula} equals A056924$(N)$, the number of strictly inferior divisors of $N$ (divisors $d$ with $d < N/d$):
\[
c_1(N) = \mathrm{A056924}(N).
\]
\end{proposition}

\begin{proof}
The condition $N/n \geq n + 1$ is equivalent to $N/n > n$ (since $N/n$ is an integer when $n \mid N$), which is $n < N/n$, i.e., $n^2 < N$, i.e., $n$ is a strictly inferior divisor of $N$.
\end{proof}

\begin{corollary}[Pronic Bonus Structure]
\label{cor:pronic}
The sequence $a(N)$ differs from A056924$(N)$ only at pronic numbers:
\[
a(N) - \mathrm{A056924}(N) = \begin{cases} m - 1 & \text{if } N = m(m+1), \\ 0 & \text{otherwise.} \end{cases}
\]
This ``pronic bonus'' grows without bound: at $N = m(m+1)$, the bonus is $m - 1 \to \infty$.
\end{corollary}

\begin{remark}[Why This Is Not a Trivial Extension]
The pronic bonus makes $a(N)$ fundamentally different from any divisor-counting sequence. At $N = 30 = 5 \cdot 6$, the bonus is 4, giving $a(30) = 8$ versus A056924$(30) = 4$. No multiplicative or additive function of $N$ can produce this bonus structure, because it depends on whether $N$ factors as a product of consecutive integers --- a property that is not preserved by the multiplicative structure of the integers. The graph-theoretic origin (the phase transition at $n = p$ in the coding clique size formula) explains why this bonus exists and why it appears precisely at pronic numbers.
\end{remark}

\subsection{Properties}

\begin{proposition}[Lower Bound]
$a(N) \geq 1$ for all $N \geq 2$, since $(n, p) = (1, N-1)$ always satisfies $N = 1 \cdot (N-1+1) = N$.
\end{proposition}

\begin{proposition}[Pronic Spike]
If $N = m(m+1)$ is pronic, then $a(N) \geq m$. This is because:
\begin{itemize}
\item $(1, N-1)$ contributes 1 from Case 1,
\item $(m, 1), (m, 2), \ldots, (m, m-1)$ contribute $m - 1$ from Case 2,
\item and possibly more from Case 1 if $m$ has divisors $\leq m$.
\end{itemize}
\end{proposition}

\begin{proposition}[Prime Values]
If $N$ is prime and not pronic, then $a(N) = 1$ (only the pair $(1, N-1)$ works).
\end{proposition}

\section{Computational Results}

\subsection{Values of $a(N)$ for $N = 2$ to $30$}

\begin{center}
\begin{tabular}{cc|cc|cc|cc}
\toprule
$N$ & $a(N)$ & $N$ & $a(N)$ & $N$ & $a(N)$ & $N$ & $a(N)$ \\
\midrule
2 & 1 & 10 & 2 & 18 & 3 & 26 & 2 \\
3 & 1 & 11 & 1 & 19 & 1 & 27 & 2 \\
4 & 1 & 12 & 5 & 20 & 6 & 28 & 3 \\
5 & 1 & 13 & 1 & 21 & 2 & 29 & 1 \\
6 & 3 & 14 & 2 & 22 & 2 & 30 & 8 \\
7 & 1 & 15 & 2 & 23 & 1 & & \\
8 & 2 & 16 & 2 & 24 & 4 & & \\
9 & 1 & 17 & 1 & 25 & 1 & & \\
\bottomrule
\end{tabular}
\end{center}

\subsection{Extended Values for $N = 31$ to $60$}

\begin{center}
\begin{tabular}{cc|cc|cc|cc}
\toprule
$N$ & $a(N)$ & $N$ & $a(N)$ & $N$ & $a(N)$ & $N$ & $a(N)$ \\
\midrule
31 & 1 & 39 & 2 & 47 & 1 & 55 & 2 \\
32 & 2 & 40 & 5 & 48 & 4 & 56 & 5 \\
33 & 2 & 41 & 1 & 49 & 2 & 57 & 3 \\
34 & 2 & 42 & 9 & 50 & 3 & 58 & 2 \\
35 & 2 & 43 & 1 & 51 & 2 & 59 & 1 \\
36 & 4 & 44 & 3 & 52 & 3 & 60 & 8 \\
37 & 1 & 45 & 3 & 53 & 1 & & \\
38 & 2 & 46 & 2 & 54 & 4 & & \\
\bottomrule
\end{tabular}
\end{center}

\subsection{Notable Values}

\begin{center}
\begin{tabular}{ccl}
\toprule
$N$ & $a(N)$ & Reason \\
\midrule
6 & 3 & Pronic $2 \cdot 3$: bonus $+1$ \\
12 & 5 & Pronic $3 \cdot 4$: bonus $+2$ \\
20 & 6 & Pronic $4 \cdot 5$: bonus $+3$ \\
30 & 8 & Pronic $5 \cdot 6$: bonus $+4$ \\
42 & 9 & Pronic $6 \cdot 7$: bonus $+5$ \\
56 & 5 & Pronic $7 \cdot 8$: bonus $+6$ \\
\bottomrule
\end{tabular}
\end{center}

\section{Derivation of Specific Values}

We illustrate the counting with three examples.

\textbf{Example 1: $N = 6$ (smallest pronic number with bonus).}

Case 1 ($n \leq p$): Divisors of 6 with $N/n \geq n+1$:
\begin{itemize}
\item $n = 1$: $p = 5 \geq 1$. Valid pair $(1, 5)$.
\item $n = 2$: $p = 2 \geq 2$. Valid pair $(2, 2)$.
\item $n = 3$: $p = 1 < 3$. Invalid.
\end{itemize}
Case 1 gives 2 pairs.

Case 2 ($n > p$): $6 = 2 \cdot 3$, so $m = 2$. Pairs: $(2, 1)$. That is 1 pair.

Total: $a(6) = 2 + 1 = 3$.

The three non-isomorphic CC-images are:
\begin{enumerate}
\item $(1, 5)$: A single cluster $K_6$ (complete graph on 6 vertices).
\item $(2, 2)$: Two clusters of $K_3$ with 2 inter-cluster edges per pair.
\item $(2, 1)$: Two clusters of $K_3$ with 1 inter-cluster edge per pair.
\end{enumerate}

\textbf{Example 2: $N = 12$ (pronic with larger bonus).}

Case 1: Divisors of 12 with $N/n \geq n+1$:
\begin{itemize}
\item $n = 1$: $p = 11$. Valid.
\item $n = 2$: $p = 5$. Valid.
\item $n = 3$: $p = 3$. Valid.
\item $n = 4$: $p = 2 < 4$. Invalid.
\end{itemize}
Case 1 gives 3 pairs.

Case 2: $12 = 3 \cdot 4$, so $m = 3$. Pairs: $(3, 1), (3, 2)$. That is 2 pairs.

Total: $a(12) = 3 + 2 = 5$.

\textbf{Example 3: $N = 7$ (prime, no bonus).}

Case 1: Only $n = 1$ divides 7. $p = 6 \geq 1$. Valid pair $(1, 6)$.

Case 2: 7 is not pronic.

Total: $a(7) = 1$.

\section{OEIS Cross-References}

\begin{itemize}
\item \textbf{A056924}: Number of strictly inferior divisors of $n$. Our sequence $a(N) = \mathrm{A056924}(N) + c_2(N)$, where $c_2(N)$ is the pronic bonus.
\item \textbf{A038548}: Number of divisors $d$ of $n$ with $d \leq \sqrt{n}$. Related via $\mathrm{A056924}(N) = \mathrm{A038548}(N) - [N \text{ is a perfect square}]$.
\item \textbf{A002378}: Oblong (pronic) numbers $m(m+1)$. These are precisely the values of $N$ where the pronic bonus is nonzero.
\item \textbf{A000040}: Prime numbers. At prime $N$ (that are not pronic), $a(N) = 1$.
\item \textbf{A006125}: Number of labeled graphs on $n$ vertices $= 2^{\binom{n}{2}}$. The CC-images are a special subclass of these.
\item \textbf{A000217}: Triangular numbers. The row $n = 1$ of the CC-image parameter space gives $N = p + 1$ for $p \geq 1$.
\end{itemize}

\section{Algorithmic Implementation}

\begin{verbatim}
def a(N):
    """Compute the CC-Image Vertex Census a(N)."""
    if N < 2:
        return 0
    # Case 1: n <= p, N = n*(p+1), p = N/n - 1 >= n
    count = 0
    for n in range(1, N + 1):
        if N % n == 0:
            p = N // n - 1
            if p >= 1 and p >= n:
                count += 1
    # Case 2: n > p, N = n*(n+1), p = 1,...,n-1
    # Check if N is pronic: N = m*(m+1) for some m
    import math
    m = int((-1 + math.isqrt(1 + 4 * N)) // 2)
    if m * (m + 1) == N and m >= 2:
        count += (m - 1)
    return count

# Compute first 59 terms (N = 2 to 60)
print([a(N) for N in range(2, 61)])
\end{verbatim}

\section{Theoretical Significance}

\begin{enumerate}
\item \textbf{Address Space of the CC-Transformation}: The sequence $a(N)$ measures how many distinct positions in the universe of simple graphs the CC-transformation occupies at each vertex count $N$. The pronic spikes reveal that the encoding is most ``expressive'' at vertex counts of the form $m(m+1)$, where the phase transition in the coding clique size formula allows multiple valid $(n, p)$ combinations.

\item \textbf{Non-Triviality of the Pronic Bonus}: The pronic bonus $m - 1$ at $N = m(m+1)$ is not an artifact of the counting method --- it reflects a genuine structural feature of the CC-transformation. When $N = m(m+1)$, the vertex count can be achieved either by $m$ clusters of size $m + 1$ (with various $p$ values in the non-dominant regime) or by $n$ clusters of size $p + 1$ (with $n \leq p$). This duality is absent at non-pronic $N$.

\item \textbf{Rado Graph Implications}: Since the Rado graph contains all countable simple graphs as induced subgraphs, the sequence $a(N)$ bounds the number of distinct CC-image types that must be accommodated at each size level within the Rado graph.
\end{enumerate}

\section{Conclusion}

The CC-Image Vertex Census sequence $a(N)$ provides a novel measure of the expressiveness of the Clique-Cluster Transformation at each vertex count. Its pronic bonus structure distinguishes it from all existing divisor-counting sequences in the OEIS and reflects a deep connection between the number-theoretic properties of $N$ and the graph-theoretic structure of CC-images.

\begin{thebibliography}{9}
\bibitem{zarzuelo2025} A. Zarzuelo Urdiales, ``Canonical Embedding of Universal $p$-Graphs, Commensurable Weighted Graphs, and Line Structures into the Rado Graph via Clique-Size Dominance Transformation,'' December 2025.
\bibitem{oeis} N. J. A. Sloane et al., The On-Line Encyclopedia of Integer Sequences, \url{https://oeis.org}.
\bibitem{hardy1979} G. H. Hardy and E. M. Wright, \textit{An Introduction to the Theory of Numbers}, 5th ed., Oxford University Press, 1979.
\end{thebibliography}

\end{document}
